\documentclass[11pt]{article}
\usepackage[margin=1in]{geometry}
\usepackage{amsmath,amssymb,microtype}
\usepackage[colorlinks=true,urlcolor=blue]{hyperref}
\title{Literature status: frequent disjoint hypercyclicity need not be dense}
\author{Automated mathematics run \texttt{fa\_banach\_001}, lane 2}
\date{11 August 2026}
\begin{document}
\maketitle

\noindent\textbf{Status:} literature already answered; both clauses of
Question 2.6 have negative answers.

\section*{Source question}
Oezguer Martin and Rebecca Sanders,
\emph{Extending Families of Disjoint Hypercyclic Operators},
arXiv:2312.07054, Question 2.6 on physical PDF page 6, ask:
\begin{quote}
Are frequently $d$-hypercyclic operators also densely $d$-hypercyclic? Do
frequently $d$-hypercyclic operators satisfy the Disjoint Blow-up/Collapse
property?
\end{quote}

\section*{Supporting answer}
Frederic Bayart, Sophie Grivaux, Etienne Matheron, and Quentin Menet,
\emph{Hereditarily Frequently Hypercyclic Operators and Disjoint Frequent
Hypercyclicity}, arXiv:2409.07103.  Section 7 restates the first clause as
Question 7.1 and says that the next theorem solves it.  Theorem 7.2, on
physical page 28, states that there are a Banach space $X$ and
$T_1,T_2\in\mathfrak L(X)$ such that $(T_1,T_2)$ is $d$-frequently
hypercyclic but not densely $d$-hypercyclic.  Its proof is completed on
physical page 32.  Hence the first clause is answered negatively.

The same example answers the second clause.  As recorded in the source
paper's implication diagram immediately before Question 2.6, the Disjoint
Blow-up/Collapse property implies dense $d$-hypercyclicity.  Therefore the
pair in Theorem 7.2 cannot have the Disjoint Blow-up/Collapse property.

\section*{Provenance and scope}
This is an exact literature resolution, not a new counterexample from the
run.  The supporting paper explicitly formulates the same density question,
labels Theorem 7.2 as its solution, and cites the source preprint
arXiv:2312.07054 as reference [39].  The negative answer to the
Blow-up/Collapse clause is the immediate contrapositive of the implication
already stated by the source.

Thus both clauses of Question 2.6 are settled negatively.  The source's
Questions 1.2 and 1.3 are outside this packet because the source itself answers
them later; those signals are same-paper extraction false positives.

\section*{Evidence}
The packet contains the original arXiv PDF as \texttt{source\_paper.pdf} and
the decisive later PDF as
\texttt{supporting\_paper\_2409.07103.pdf}.  The source question is on physical
page 6.  The supporting question and theorem are on physical page 28, and the
proof concludes on physical page 32.

\begin{thebibliography}{9}
\bibitem{MS2023}
O. Martin and R. Sanders,
\emph{Extending Families of Disjoint Hypercyclic Operators},
arXiv:2312.07054 (2023).

\bibitem{BGMM2024}
F. Bayart, S. Grivaux, E. Matheron, and Q. Menet,
\emph{Hereditarily Frequently Hypercyclic Operators and Disjoint Frequent
Hypercyclicity}, arXiv:2409.07103 (2024).
\end{thebibliography}

\end{document}
